\documentclass[../BTL.tex]{subfiles}
\begin{document}

% ─────────────────────────────────────────────────────────
\section{Kết quả đạt được}
\label{section:4.1}

Sau quá trình nghiên cứu và thực hiện, đồ án đã đạt được đầy đủ các mục tiêu
đề ra tại mục~\ref{section:1.2}.

\subsection{Về lý thuyết}
\label{section:4.1.1}

\begin{itemize}[itemsep=0.4em]
    \item Nắm vững nền tảng lý thuyết về hệ thống gợi ý, phân biệt rõ ba
    phương pháp tiếp cận: \textit{Content-based Filtering},
    \textit{Collaborative Filtering} và \textit{Hybrid}.

    \item Hiểu sâu cơ chế hoạt động của \textit{User-based CF} và
    \textit{Item-based CF}, bao gồm \textit{mean-centered cosine similarity}
    và \textit{adjusted cosine similarity}.

    \item Nắm vững thuật toán \textit{Matrix Factorization} (Funk SVD), hàm
    mất mát với \textit{regularization} L2 và cơ chế SGD cập nhật tham số.

    \item Hiểu và áp dụng đúng bốn chỉ số đánh giá: MAE, RMSE,
    Precision@$K$, Recall@$K$; phân biệt khi nào dùng chỉ số nào.
\end{itemize}

\subsection{Về kỹ thuật}
\label{section:4.1.2}

\begin{itemize}[itemsep=0.4em]
    \item Triển khai thành công ba mô hình CF \textit{from scratch} bằng
    Python/NumPy, không phụ thuộc thư viện ML chuyên dụng
    (\texttt{Scikit-learn} chỉ được dùng để tính chỉ số MAE/RMSE).

    \item Áp dụng \textit{vectorized computation} (nhân ma trận NumPy) để
    tính ma trận tương đồng hiệu quả, thay thế vòng lặp Python thuần túy.

    \item Cài đặt \textit{temporal split} tránh \textit{data leakage}.
    Thực hiện \textit{grid search} tối ưu siêu tham số ($K$, $\alpha$,
    $\lambda$).

    \item Xây dựng hệ thống \textit{backend/frontend} hoàn chỉnh: FastAPI
    RESTful API + giao diện web \textit{responsive dark theme}.
\end{itemize}

\subsection{Về kết quả thực nghiệm}
\label{section:4.1.3}

Bảng~\ref{tab:summary} tổng hợp kết quả thực nghiệm trên bộ dữ liệu
\texttt{MovieLens ml-100k} ($100.000$ \textit{ratings}, $943$ users,
$1.682$ phim).

\begin{table}[H]
    \centering
    \caption{Tổng hợp kết quả thực nghiệm trên MovieLens ml-100k}
    \label{tab:summary}
    \begin{tabular}{lllllll}
        \toprule
        \textbf{Mô hình}
            & \textbf{MAE} & \textbf{RMSE}
            & \textbf{P@10} & \textbf{R@10}
            & \textbf{Coverage} & \textbf{Xếp hạng} \\
        \midrule
        Baseline (Mean)
            & ${\approx}0{,}945$ & ${\approx}1{,}126$
            & --- & ---
            & $100\%$ & --- \\
        User-based CF ($K=30$)
            & $0{,}7746$ & $0{,}9863$
            & $0{,}0047$ & $0{,}0034$
            & ${\approx}98\%$ & 2 \\
        Item-based CF ($K=30$)
            & $0{,}7996$ & $1{,}0211$
            & $0{,}0152$ & $0{,}0092$
            & ${\approx}97\%$ & 3 \\
        \textbf{Matrix Factorization}
            & $\mathbf{0{,}7706}$ & $\mathbf{0{,}9758}$
            & $\mathbf{0{,}0518}$ & $\mathbf{0{,}0484}$
            & $\mathbf{100\%}$ & \textbf{1} \\
        \bottomrule
    \end{tabular}
\end{table}

\textit{Matrix Factorization} (Funk SVD, $K=20$) đạt kết quả tốt nhất trên
tất cả các chỉ số: RMSE = $0{,}9758$ (cải thiện $13{,}3\%$ so với
\textit{baseline}), Precision@10 = $0{,}0518$ (gấp ${\approx}11$ lần
\textit{User-based CF}) và \textit{coverage} $100\%$. \textit{User-based CF}
đứng thứ hai với RMSE = $0{,}9863$, trong khi \textit{Item-based CF} bị ảnh
hưởng nhiều hơn bởi độ thưa thớt $93{,}7\%$ của dữ liệu.

% ─────────────────────────────────────────────────────────
\section{Hạn chế của đề tài}
\label{section:4.2}

Bên cạnh những kết quả đạt được, đề tài còn tồn tại một số hạn chế:

\begin{itemize}[itemsep=0.5em]
    \item \textbf{Cold start problem} (vấn đề khởi động nguội): Các mô hình
    CF và MF không thể gợi ý cho người dùng mới (chưa có
    \textit{rating} nào) hoặc phim mới (chưa được ai đánh giá). Đây là
    hạn chế cơ bản của \textit{Collaborative Filtering}. Hệ thống hiện xử
    lý bằng cách \textit{fallback} về \textit{rating} trung bình,
    nhưng chưa giải quyết triệt để.

    \item \textbf{Scalability} (khả năng mở rộng): Ma trận tương đồng
    $943 \times 943$ (UBCF) và $1.615 \times 1.615$ (IBCF) còn quản lý
    được trong RAM. Tuy nhiên với dataset lớn hơn (hàng triệu users/items),
    cần các giải pháp \textit{Approximate Nearest Neighbors} (ANN) như
    Faiss hay HNSW.

    \item \textbf{SGD chưa tối ưu hóa}: Vòng lặp SGD thuần Python
    (${\approx}80.000$ iterations/epoch) còn chậm do không được vector hóa
    hoàn toàn. Có thể tăng tốc đáng kể bằng Cython hoặc PyTorch.

    \item \textbf{Chỉ dùng \textit{rating} tường minh}: Mô hình chưa khai
    thác dữ liệu ngầm định (\textit{implicit feedback}) như lượt xem, thời
    gian xem, lịch sử tìm kiếm — vốn phong phú hơn nhiều trong thực tế.

    \item \textbf{Chưa có Content-based}: Thông tin nội dung phim (thể loại,
    đạo diễn, diễn viên) chưa được tận dụng để bổ sung đặc trưng cho mô
    hình hay giải quyết \textit{cold start}.
\end{itemize}

% ─────────────────────────────────────────────────────────
\section{Hướng phát triển tương lai}
\label{section:4.3}

Dựa trên những kết quả và hạn chế trên, các hướng phát triển tiếp theo được
đề xuất:

\begin{enumerate}[itemsep=0.5em]
    \item \textbf{Neural Collaborative Filtering (NCF)}: Thay thế tích vô
    hướng $p_u \cdot q_i$ bằng mạng neural sâu để học được tương tác phi
    tuyến phức tạp giữa user và item. NCF đã được chứng minh vượt trội so
    với MF truyền thống trên nhiều \textit{benchmark}~\cite{he2017neural}.

    \item \textbf{Hybrid Recommender System}: Kết hợp \textit{Collaborative
    Filtering} với \textit{Content-based Filtering} (thông tin thể loại, đạo
    diễn, diễn viên) để giải quyết \textit{cold start problem} và cải thiện
    độ đa dạng gợi ý.

    \item \textbf{Implicit Feedback Model}: Áp dụng các mô hình như BPR
    (\textit{Bayesian Personalized Ranking}) hoặc ALS
    (\textit{Alternating Least Squares}) trên dữ liệu ngầm định thay vì
    chỉ \textit{rating} tường minh.

    \item \textbf{Đánh giá đa chiều}: Bổ sung thêm các chỉ số
    NDCG@$K$ (\textit{Normalized Discounted Cumulative Gain}),
    MRR (\textit{Mean Reciprocal Rank}) và \textit{Coverage} để đánh giá
    toàn diện hơn chất lượng danh sách gợi ý.

    \item \textbf{Triển khai production}: Tích hợp \textit{caching} (Redis),
    \textit{authentication}, \textit{logging} và \textit{containerization}
    (Docker) để triển khai hệ thống trong môi trường thực tế.
\end{enumerate}

\end{document}
